\Qhead{4.1}{Принцип работы, устройство и параметры лазеров}{Принцип работы, устройство и параметры лазеров (примеры: гелий-неоновый лазер, лазер на рубине, полупроводниковый лазер).}
\Subfile{4.1.1}{sections/04_generation/q01/s01}

\recap{Накачка создаёт инверсию, среда усиливает поле, резонатор выбирает моды, насыщение ограничивает рост. Три примера различаются способом накачки, уровневой схемой и механизмом ограничения носителей.}
\board{Четырёхуровневая схема и порог; сравнение He--Ne, рубина и инжекционного лазера.}
\oral{Почему резонансно накачиваемая замкнутая двухуровневая система не создаёт устойчивой инверсии?}
\bait{Можно оставить высокий порог рубина: приходится опустошать населённый основной уровень, в отличие от эффективной четырёхуровневой схемы.}
\sources{Университетские слайды «Основы оптоэлектроники»~\cite{optoelectronics}, раздел источников; Kärtner~\cite{kartnerfull}, главы~2 и 4. Скоростная модель приведена в согласованной нормировке.}
